\section{Conclusions}
\label{chapter:05}

% block-id: 05.conclusions.p1
\hypertarget{05.conclusions.p1}{}
Under common phasor, length, and area definitions, this paper establishes verification relationships among round-wire analysis, two-dimensional common-terminal coupling, and finite-length longitudinal PEEC. The finest two-dimensional internal-field mesh has a relative analytical resistance error of 1.79×10⁻⁶. For parallel-connected strands, uniform currents are characterized by a zero differential-subspace projection of coupled impedance and imposed-field induction. Complex currents in the 640-strand straight bundle show that common terminals and fine subdivision must be assessed together with magnetic coupling. Longitudinal PEEC meets the relevant sensitivity targets for the tested lengths and integration settings; the final 5024-degree-of-freedom resistance differs by 0.044582\% from its equal-area Bessel reference.

% block-id: 05.conclusions.p2
\hypertarget{05.conclusions.p2}{}
These results provide local-field references, common-terminal constraints, and finite-length computational benchmarks for research on actual transposition. RT0 imposed-field cases, CPU/GPU comparisons, and three-dimensional diagnostics further describe the implementations' ranges. Ranking actual Litz structures requires three-dimensional verification with end regions and transverse degrees of freedom, followed by thermal assessment and controlled experiments.

% block-id: 05.current.evidence
\hypertarget{05.current.evidence}{}
The current-device environment migration reproduced the analytical calibration and the small Maxwell port-and-power case. These runs verify the migrated execution path without enlarging the convergence scope of the retained Maxwell mesh or changing the failed A–φ gauge status. Recently checked semi-analytical and winding-loss literature supplies candidate methods for subsequent work. Their assumptions about winding fields and homogenization require explicit comparison with this paper's isolated-wire and common-terminal models before quantitative results are transferred. \cite{C01,ZHU2026}
